% !PRINTED 217
\weberpage{217}

%<<CONT
und die Relationen
\weq{6}{\Theta\chi = \chi^2\psi\Theta^3, \quad \Theta^2\chi = \chi\psi, \quad \Theta^3\chi = \chi^2\Theta \\
\psi\chi = \chi\psi\Theta^2, \quad \psi\chi^2 = \chi^2\Theta^2.}
Diese Relationen in Verbindung mit (2) zeigen, dass die
Substitutionen
\weq{7}{\chi^{\lambda}\psi^{\mu}\Theta^{\nu}, \quad \lambda = 0, 1, 2; \quad \mu = 0, 1; \quad \nu = 0, 1, 2, 3}
in der That eine Gruppe bilden, weil man mit ihrer Hülfe jede
Substitution der Form
\[
\chi^{\lambda}\psi^{\mu}\Theta^{\nu}\chi, \quad \chi^{\lambda}\psi^{\mu}\Theta^{\nu}\psi,
\]
und folglich durch Wiederholung auch jede Substitution
\[
\chi^{\lambda}\psi^{\mu}\Theta^{\nu}\,\chi^{\lambda'}\psi^{\mu'}\Theta^{\nu'}
\]
in die Form (7) bringen kann. Die Gruppe kann explicite in
der Form
\weq{8}{i^{\nu}x, \quad \frac{i^{\nu}}{x}, \quad i^{\nu}\,\frac{x - i^{\nu'}}{x + i^{\nu'}}, \quad \nu, \nu' = 0, 1, 2, 3}
dargestellt werden und ist vom 24\textsuperscript{sten} Grade. Es ist die
\gesperrt{Octaëdergruppe}.

Die Gruppe hat einen Normaltheiler 12\textsuperscript{ten} Grades, den man
in der Form $\chi^{\lambda}\psi^{\mu}\Theta^{2\nu}$ darstellen kann und der eine
Tetraëdergruppe ist.

Die Darstellung wird in gewisser Beziehung einfacher, wenn
man an Stelle von $\psi$ ein neues Element
\weq{9}{\omega = \psi\Theta = \frac{-i}{x}}
einführt, was ebenso wie $\psi$ von der zweiten Ordnung ist. Dann
kann die Gruppe dargestellt werden durch
\[
\chi^{\lambda}\omega^{\mu}\Theta^{\nu},
\]
und an Stelle der Relationen (2) und (6) treten die folgenden:
\[
\omega\chi = \chi^2\omega, \quad \omega\chi^2 = \chi\omega
\]
\weq{10}{\Theta\chi = \chi^2\omega\Theta^2, \; \Theta^2\chi = \chi\omega\Theta^3, \; \Theta^3\chi = \chi^2\Theta, \; \Theta^2\chi^2 = \chi^2\omega\Theta \\
\Theta\omega = \omega\Theta^3, \quad \Theta^2\omega = \omega\Theta^2, \quad \Theta^3\omega = \omega\Theta.}
Alle diese Relationen aber ergeben sich als Folgerungen
aus den vieren:
\weq{11}{\omega\chi = \chi^2\omega, \quad \Theta\omega = \omega\Theta^3, \quad \Theta\chi = \chi^2\omega\Theta^2, \quad \Theta^2\chi = \chi\omega\Theta^3,}
in Verbindung mit den die Grade ausdrückenden Formeln:
\weq{12}{\chi^3 = 1, \quad \omega^2 = 1, \quad \Theta^4 = 1.}
